% page 34 of the facsimile (image p-033 of the scan)
% block 165.1 x 229.0 mm ; left margin 8.0 mm
\pgc{-12.700mm}{0.000mm}{
\l{\cel{7}26}
\l{}
\l{\sou{anaerobia}.\cel{2}(\sou{biol}.) (\sou{Pri la organismi mikroskopala}). Qua povas}
\l{\cel{3}vivar en medio sen-oxa e sen-aera. - DEFIS.}
\l{}
\l{\sou{anafazo}. (\sou{biol}.)Fazo di la konstituco di la nuklei filio lor}
\l{\cel{3}la divido mitosala di la nuklei, divido per qua finas mi-}
\l{\cel{3}toso. - DEFIS.}
\l{}
\l{\sou{anagalo}. (\sou{bot.})\cel{2}Planto singla-yara, ek la familio \textquotedbl{}primulacei\textquotedbl{}.}
\l{\cel{3}L.\cel{1}\sou{anagallis}. DEFIS.}
\l{}
\l{\sou{anagramo}.\cel{2}Vorto formacita ek la transpozo di la literi dialtra}
\l{\cel{3}vorto. - DEFIS.}
\l{}
\l{\sou{anakoluto}.\cel{2}Elipso en qua onu omisas, en frazo, la korelativo}
\l{\cel{3}gramatikala di vorto expresata. (Ex.: \textquotedbl{}\sou{Qua}\cel{1}servas bone sua}
\l{\cel{3}patrio / Ne bezonas ancestri\textquotedbl{} (Voltaire). La omiso di\cel{1}\sou{ta}\cel{1}avan}
\l{\cel{3}\sou{qua}\cel{1}esas anakoluto). - DEFIRS.}
\l{}
\l{\sou{anakoreto}. Religiano qua vivas sole. - DEFIRS.}
\l{}
\l{\sou{anakronismo}.\cel{2}Eroro qua konsistas ye konsiderar ke la fakto}
\l{\cel{3}eventis ante sua dato eventala o (per extenso di la senco),}
\l{\cel{3}ye dato altra kam la lua. - DEFIRS.}
\l{}
\l{\sou{analo}. Singla de la naracuri pri la eventaji di la yaro koncern-}
\l{\cel{3}ata. - DEFIS.}
\l{}
\l{\sou{analgeziar}.\cel{2}(\sou{trans}.) Produktar anestezio partala di la pelo o di}
\l{\cel{3}la mukozi. - EFIS.}
\l{}
\l{\sou{analitika}. I. Qua procedas per analizado. - II.\cel{1}\sou{Geometrio ana}-}
\l{\cel{3}\sou{litika}\cel{1}: Apliko di algebro a geometrio, qua, retroduktante}
\l{\cel{3}la nociono \textquotedbl{}figuro\textquotedbl{} ad olta di \textquotedbl{}situeso\textquotedbl{}, ed olta di \textquotedbl{}situ-}
\l{\cel{3}eso\textquotedbl{} ad olta di \textquotedbl{}dimensiono\textquotedbl{}, reprezentas la loki geometriala}
\l{\cel{3}per equacioni, o deduktas, de la analizo di la equacioni, la}
\l{\cel{3}karakterizivi di la figuri. - III.\cel{1}\sou{Linguo analitika}\cel{1}: Qua}
\l{\cel{3}tendencas deskompozar la moyeni per qui onu expresas la pens-}
\l{\cel{3}aji. - DEFIRS.}
\l{}
\l{\sou{analizar}.\cel{2}(\sou{trans}.) Dividar ensemblo ad olua parti. - DEFIRS.}
\l{}
\l{\sou{analoga}. Qua havas kelka karakteri komuna kun altra kozo.-DEFIRS.}
\l{}
\l{\sou{anamorfoso}. (\sou{Mat.})\cel{2}Fenomeno fizikala qua eventas kande la}
\l{\cel{3}desegnuro deformita de objekto, pozate perpendikle ye la axo}
\l{\cel{3}di spegulo cilindra o kona, donas, per reflekto, imajo de}
\l{\cel{3}la objekto sen deformo. - DEFIRS.}
\l{}
\l{\sou{ananaso}.(\sou{bot.)}\cel{2}Frukto saporoza e parfumoza, di planto dornoza,}
\l{\cel{3}ek la familio \textquotedbl{}bromeliacei\textquotedbl{}, di qua la formo esas ovatra.}
\l{\cel{3}L.\cel{1}\sou{Ananasa sativa}. - DFIRS.}
}
